
\chapter{نتیجه‌گیری و کارهای آینده}
\label{ch:conclusions}

\section{جمع‌بندی}
\label{sec:summary}

هدف این پروژه، طراحی، پیاده‌سازی و ارزیابی یک نمونه اولیه مدل زبانی بزرگ
گفتار‌به‌گفتار تمام‌دوطرفه برای زبان فارسی بود. برای کاهش ریسک اجرایی، دو
مسیر به‌طور موازی پیش برده شد: تطبیق کم‌پارامتر یک مدل سرتاسری
گفتار‌به‌گفتار متن‌باز، و یک سامانه آبشاری ماژولار با کنترل
تمام‌دوطرفه بیرونی. نتیجه، نامتقارن بود.

\مهم{مسیر آبشاری} به یک سامانه کارکردی و قابل ارزیابی رسید. در یک پنل
ثابت و ازپیش‌قفل‌شده اعتبارسنجی، زنجیره با پاسخ‌گوی \en{Qwen2.5-0.5B} در هر
\N{9} ردیف رونویسی فارسی، پاسخ فارسی و صوت غیرساکت تولید کرد و هیچ‌گاه به
پاسخ جایگزین قاعده‌محور بازنگشت. پس از عبور از یک مرحله صلاحیت هفت‌ردیفی،
آزمون نهایی \N{40} ردیفی گروه‌ناهمپوشان روی پاسخ‌گوی \en{Qwen3-4B} از هر هفت
شرط خودکار ازپیش‌تعریف‌شده عبور کرد: میانگین ارتباط پاسخ از \N{1.425} به
\N{3.150}، میانگین انسجام به \N{3.325}، نرخ برد زوجی به \N{67.5} درصد رسید
و هر \N{240} فراخوانی داور معتبر بود. این سنجه داوری هم‌خانواده است و
جایگزین ارزیابی انسانی نیست. سنجه خودکار بازبازشناسی روی همان \N{9} خروجی
\en{0.5B}، نرخ خطای نویسه \N{7.14} درصد را اندازه گرفت.

\مهم{مسیر مستقیم} در شش نسل آزمایشی به هدف نرسید. در هر نسل، زیان تحت
تحمیل معلم کاهش یافت اما شرط اجرای خودرگرسیو نه‌ردیفی هیچ‌گاه برآورده نشد
و هیچ چک‌پوینتی واجد شرایط استقرار نشد. آخرین نسل با کاهش
\N{32.20} درصدی زیان متنی درون‌نمونه نشان داد که آداپتر ظرفیت یادگیری
هدف را دارد، اما هر چهار چک‌پوینت آن \N{0} از \N{9} گرفتند. این مسیر
به‌عنوان یک نتیجه منفی مستند با سیگنال یادگیری مثبت گزارش شده است.

\مهم{مجموعه داده} از یک موجودی \N{775.887} ساعتی به \N{6,754} جفت
«گفتار کاربر و پاسخ» با \N{108.584} ساعت صوت استریوی ممیزی‌شده رسید که
درون بازه \N{100} تا \N{200} ساعت مورد نظر تعریف پروژه قرار دارد، با
تفکیک گروهی بر پایه نشست منبع و نشت گروهی صفر.

\مهم{ماژول تشخیص وقفه} با ویژگی‌های انرژی، زیر‌و‌بمی و \en{MFCC} و یک
طبقه‌بند تقویت گرادیانی پیاده شد و روی رخدادهای صوت واقعی با نشست‌های
کنارگذاشته به دقت نقطه‌ای \N{81.06} درصد رسید؛ کران پایین فاصله اطمینان
بلوکی در سطح نشست \N{74.44} درصد است و برچسب‌ها خودکارند، بنابراین عبور
رسمی از \N{80} درصد اثبات نشده است. خط پایه انرژی دقت \N{78.03} درصد گرفت
و نرخ رد نادرست آن \N{43.94} درصد بود.

برای چهار نیازمندی هنوز شواهد رسمی کامل گردآوری نشد: تأخیر سرتاسری
\N{500} میلی‌ثانیه، عبور رسمی تشخیص وقفه از \N{80} درصد با برچسب مستقل
انسانی، ارزیابی روی پردازنده گرافیکی فیزیکی \N{12} تا \N{24} گیگابایتی، و
مطالعه انسانی با \N{5} تا \N{10} گویشور فارسی. ازاین‌رو، این نوشتار
درباره تعامل سالمندان فارسی‌زبان ادعای تجربی مطرح نمی‌کند.

\section{پاسخ به پرسش پژوهش}
\label{sec:answer}

پرسش عملی این پروژه، که در بخش~\ref{sec:problem} مطرح شد، دو بخش داشت.

بخش نخست: با بودجه محاسباتی و زمانی یک پروژه کارشناسی، تا چه اندازه
می‌توان یک سامانه گفتار‌به‌گفتار فارسی تمام‌دوطرفه ساخت؟ پاسخ این است که
یک سامانه «کارکردی» با معماری ماژولار و کنترل وقفه بیرونی قابل ساخت
است و مرتبط‌بودن پاسخ آن در یک آزمون خودکار هم‌خانواده بهتر از خط پایه
دیده شد، نه کیفیت ادراکی یا مستقل؛ اما یک مدل سرتاسری فارسی، با
آموزش‌دهنده‌های رسمی موجود و در این بودجه، قابل دستیابی نبود.

بخش دوم، که مهم‌تر بود: چه چیزی را می‌توان با شواهد قابل ممیزی
«اثبات» کرد؟ پاسخ در جدول~\ref{tab:req-status} خلاصه شده است. دو
نیازمندی ن۱ و ن۴ شواهد کافی دارند؛ دو نیازمندی ن۲ و ن۵ پیاده‌سازی شده‌اند
اما شواهد رسمی آن‌ها کامل نیست؛ و برای سه نیازمندی ن۳، ن۶ و ن۷ شواهد
رسمی موجود نیست.
سهم روشمند این نوشتار در همین تفکیک است: پروژه به‌جای ارائه یک ادعای
کلی موفقیت، برای هر سنجه مرز ادعای مجاز آن را مشخص کرده و در کد نیز یک
سازوکار شکست‌بسته مانع از ارتقای تقریب‌ها به شواهد رسمی شده است.

\section{کارهای آینده}
\label{sec:future}

پیشنهادهای کار آینده به‌ترتیب اولویت و به تفکیک «تکمیل شواهد» و
«بهبود فنی» ارائه می‌شود.

\subsection{تکمیل شواهد باقی‌مانده}

این چهار مورد به پژوهش تازه نیاز ندارند و تنها اجرای پروتکل‌های
ازپیش‌نوشته‌شده را لازم دارند.

\begin{enumerate}
\item \مهم{اندازه‌گیری تأخیر روی سخت‌افزار هدف.} اجرای پروتکل ثبت‌شده روی
      یک پردازنده گرافیکی فیزیکی در بازه \N{12} تا \N{24} گیگابایت، با
      یک مرورگر واقعی و میکروفون فیزیکی، و گزارش
      $T_{\text{first\_audio}}$ و $T_{\text{barge\_in}}$ بر پایه رویداد
      تأییدشده کارخواه. این کار نیازمندی ن۳ و ن۶ را به کلاس شواهد رسمی
      می‌برد.

\item \مهم{مطالعه انسانی.} اجرای مطالعه با \N{5} تا \N{10} گویشور فارسی
      شامل دست‌کم دو نفر بالای \N{60} سال، با سند رضایت‌نامه و برگه
      امتیازدهی موجود و در حالت نگه‌داری فقط‌ویژگی. این کار نیازمندی ن۷ و
      قلم چهارم خروجی پروژه را تحویل می‌دهد.

\item \مهم{برچسب‌گذاری مستقل رخدادهای وقفه.} بازبینی شنیداری نمونه‌ای از
      رخدادهای آشکارساز توسط دو یا چند بازبین انسانی مستقل، همراه با
      گزارش توافق میان‌بازبینی، و سپس ارزیابی مجدد آشکارساز روی برچسب
      انسانی. این کار نیازمندی ن۵ را قابل اثبات می‌کند.

\item \مهم{بازبینی شنیداری مجموعه داده.} اجرای بازبینی چهل‌ردیفی و
      بیست‌و‌چهار‌ردیفی که تحت انصراف زمانی معلق ماند، تا بتوان بخشی از
      مجموعه داده را «تأییدشده توسط انسان» نامید.
\end{enumerate}

\subsection{بهبود تأخیر سامانه آبشاری}

مشاهده بخش~\ref{sec:res-cascade} نشان داد که زمان تولید کامل پاسخ با طول
ورودی رشد می‌کند و میانه آن چند ثانیه است. سه مداخله می‌تواند این را
به‌طور اساسی بهبود دهد.

نخست، \مهم{تولید جویباری}. در معماری کنونی، کل متن پاسخ تولید و سپس کل
موج صوتی ساخته می‌شود. اگر تولید متن به‌صورت جویباری انجام شود و هر
عبارت کامل بی‌درنگ به موتور گفتاری برود، نخستین قطعه صوت می‌تواند پیش از
تکمیل پاسخ پخش شود. این تغییر سنجه نخستین صوت را از زمان تولید کامل جدا
می‌کند و تنها راه عملی نزدیک‌شدن به هدف \N{500} میلی‌ثانیه در این معماری
است.

دوم، \مهم{کوتاه‌کردن بافت ورودی}. ردیف \N{32} پنل، با \N{5031} نویسه
رونویسی، کندترین نوبت را داشت. محدودکردن ورودی مدل زبانی به آخرین
اظهار کاربر، به‌جای کل متن رونویسی، باید این دنباله را حذف کند.

سوم، \مهم{مدل پاسخ‌گوی کوچک‌تر یا کوانتیزه.} دو چک‌پوینت جایگزین محلی در
جریان پروژه بررسی و تنها در سطح راه‌اندازی آزموده شدند. یک نسخه
کم‌حافظه‌تر می‌تواند هم تأخیر و هم تناسب با بازه حافظه هدف را بهبود دهد.
هرگونه جانشینی مدل پاسخ‌گو باید با یک آزمون معنایی تازه و
«ازپیش‌تعریف‌شده» روی ردیف‌های گروه‌ناهمپوشان انجام شود؛ پنل آزمون
نهایی کنونی یک‌بار باز شده است و بازاستفاده از آن برای انتخاب مدل مجاز
نیست.

\subsection{مسیرهای جایگزین برای مدل مستقیم}

بر پایه تحلیل بخش~\ref{sec:disc-gap}، سه راهبرد امیدبخش‌تر از تکرار
دستور کنونی به نظر می‌رسد.

نخست، \مهم{رویکرد مدل منجمد}. معماری
\en{Freeze-Omni}~\cite{wang2025freezeomni} وزن‌های مدل زبانی را منجمد نگه
می‌دارد و تنها کدگذار و کدگشای گفتاری را آموزش می‌دهد. این طراحی دقیقاً
مسئله‌ای را هدف می‌گیرد که در این پروژه شکست خورد: چون توزیع زبانی مدل
پایه دست‌نخورده می‌ماند، مشکل فروریختن رفتار تولیدی و بازگشت به زبان
پایه پیش نمی‌آید. برای فارسی این احتمالاً کوتاه‌ترین مسیر به یک مدل
سرتاسری کارکردی است.

دوم، \مهم{اصلاح کیفیت هدف آموزشی}. تحلیل بخش~\ref{sec:disc-reference}
نشان داد که نوبت‌های هم‌ترازشده خودکار، اهداف ضعیفی برای پاسخ‌دهی
دستور‌پیرو هستند. یک مرحله گزینش که تنها جفت‌هایی را نگه می‌دارد که در
آن‌ها نوبت دوم واقعاً پاسخی به نوبت اول است، حتی اگر حجم داده را چند
برابر کاهش دهد، احتمالاً هدف آموزشی بهتری می‌سازد. مجموعه داده کنونی
برای یادگیری «زمان‌بندی» گفتار و پدیده‌های تعاملی مناسب است و برای
یادگیری «محتوای» پاسخ نامناسب؛ تفکیک این دو نقش در طراحی آموزش
می‌تواند مؤثر باشد.

سوم، \مهم{تفکیک صریح مرحله زبان و مرحله گفت‌وگو}. به‌جای یک مرحله تطبیق
که هم‌زمان زبان تولیدی و رفتار گفت‌وگویی را هدف می‌گیرد، می‌توان نخست با
داده تک‌گوینده فراوان فارسی، مانند
\en{ManaTTS}~\cite{qharabagh2025manatts}، توانایی تولید گفتار فارسی را
تثبیت کرد و تنها در مرحله دوم رفتار گفت‌وگویی و نوبت‌گیری را روی داده
مکالمه آموخت.

\subsection{بهبود ماژول تشخیص وقفه}

دو بهبود مشخص قابل انجام است. نخست، \مهم{کاهش نرخ رد نادرست}. مقدار
کنونی \N{28.79} درصد است، یعنی نزدیک به یک وقفه از هر سه وقفه نادیده
می‌ماند. افزودن ویژگی‌های نوایی بلندبرد‌تر و پنجره متن طولانی‌تر
می‌تواند مؤثر باشد. دوم، \مهم{استفاده از نشانه سمت سامانه}. آشکارساز
کنونی فقط جریان میکروفون را می‌بیند. اگر سیگنال صوتی که سامانه در همان
لحظه پخش می‌کند نیز به‌عنوان ورودی مرجع در دسترس باشد، می‌توان
پس‌خورد آکوستیکی بلندگو را از گفتار واقعی کاربر جدا کرد و نرخ پذیرش
نادرست \N{9.09} درصدی را کاهش داد.

\section{سخن آخر}
\label{sec:closing}

این پروژه نشان داد که ساخت یک نمونه اولیه گفتار‌به‌گفتار فارسی با کنترل
تمام‌دوطرفه، در بودجه یک پروژه کارشناسی، امکان‌پذیر است؛ و در همان حال
نشان داد که تطبیق یک مدل سرتاسری گفتار‌به‌گفتار به یک زبان جدید، با
دستور تطبیق کم‌پارامتر و داده مشتق‌شده خودکار، به‌سادگی به دست نمی‌آید.

نتیجه منفی مسیر مستقیم، از دید ما ارزشی کمتر از نتیجه مثبت مسیر آبشاری
ندارد. این نتیجه، همراه با شواهد ثبت‌شده آن، یک هشدار روشمند برای
پژوهش‌های بعدی است: در این حوزه، کاهش زیان آفلاین می‌تواند به‌طور کامل
گمراه‌کننده باشد و آزمون تولید در زمان اجرا باید بخشی از معیار انتخاب
باشد. اگر این نوشتار تنها یک چیز را منتقل کند، همین است.
